\section{Conservation Properties and Mass Balance}
\label{sec:conservation-diagnostic}

We verify the numerical consistency of the coupled solver by monitoring the global mass balance and energy evolution.
The mass balance error $\mathcal{E}_M(t)$ is computed as the normalized difference between the rate of mass change $dM/dt$ and the integrated source term $Q(t) = \int_\Omega \mathcal{S}(\rho, \dots)\,dV$:
\begin{equation}
    \mathcal{E}_M(t) = \frac{\left| \frac{dM}{dt} - Q(t) \right|}{\max(M(t), \epsilon)},
\end{equation}
where the error is normalized by the total system mass $M(t)$ to quantify the spurious mass generation rate relative to the system size (avoiding artificial singularities at steady state where $dM/dt \to 0$).

\begin{figure}[htbp]
    \centering
    \safeincludegraphics[width=\linewidth]{conservation_diagnostic_box_physio.png}
    \caption{Conservation diagnostic for the physiological benchmark case.
    (a) Evolution of total bone mass $M(t)$ (solid) and relative change $\Delta M/M_0$ (dashed).
    (b) Mass balance check: instantaneous mass rate $dM/dt$ (solid) matches the integrated source term (dashed). The normalized balance error (red, right axis) remains negligible ($\sim 10^{-7}\%\,M/\text{day}$), confirming discrete conservation.
    (c) Strain energy evolution $W(t)$, showing smooth convergence to a homeostatic state.}
    \label{fig:conservation-physio}
\end{figure}

\begin{figure}[htbp]
    \centering
    \safeincludegraphics[width=\linewidth]{conservation_diagnostic_box_stiff.png}
    \caption{Conservation diagnostic for the stiffened benchmark case.
    The system exhibits a significant mass reduction ($\sim 17\%$) as it adapts to the stiffer material parameters.
    Despite larger transients, the mass balance error (b) remains effectively zero ($\sim 10^{-7}\%\,M/\text{day}$), demonstrating the robustness of the implicit coupling scheme even under large remodeling rates.}
    \label{fig:conservation-stiff}
\end{figure}

Comparison of the two cases reveals that while the stiffened model undergoes more aggressive remodeling (larger $|dM/dt|$ peak in panel b), the numerical scheme maintains conservation with equal precision. The steady state is reached asymptotically in both cases, with the net mass rate converging to zero and the strain energy stabilizing, consistent with the theoretical prediction of homeostatic equilibrium.
